\section[Automorphisms of Monge-Amp{\`e}re systems with Lagrangian pairs]{Automorphisms of Monge--Amp{\`e}re systems\\ with Lagrangian pairs}
\label{Automorphisms}

Let us consider the automorphism pseudo-group Aut($\mathcal{M}$) of
all local isomorphisms of a bi-decomposable Monge--Amp{\`e}re system
$\mathcal{M} = \langle \theta, d\theta, \omega \rangle$
with a Lagrangian pair $(E_1,E_2)$
on a~contact manifold $(M,D)$ of dimension $2n+1$.

By Theorems \ref{bi-decomposable-form} and~\ref{Uniqueness}, in the case $n \geq 3$,
any automorphism of any bi-decomposable Monge--Amp{\`e}re system with a~Lagrangian pair $(E_1, E_2)$ preserves the Lagrangian pair $(E_1, E_2)$
up to the interchange of~$E_1$,~$E_2$.
Therefore inf\/initesimal symmetries
of bi-decom\-po\-sable Monge--Amp{\`e}re systems
are studied based on inf\/initesimal symmetries of Lagrangian contact structures.

Now let $\overline{G}$ be
\begin{gather*}
\overline{G} = \left\{\left. \begin{pmatrix}
                       c^2 & 0 & 0 \\
                        b_1 & cA &   \mathrm{O}_n \\
                        b_2 &   \mathrm{O}_n &   c\, {}^t\!\! A^{-1}
                    \end{pmatrix}
                    \, \right\vert \,
c \in {\mathbf{R}}^{\times}, \,A\in \mathrm{SL}(n,{\mathbf{R}}),\, b_1, b_2 \in {\mathbf{R}}^n
\right\}.
\end{gather*}

\begin{Prop}
The bi-decomposable Monge--Amp{\`e}re systems with Lagrangian pairs
are in bijective correspondence with $\overline{G}$-structures of type~$\mathfrak{m}$.
\end{Prop}

\begin{proof}
We consider, as in Section~\ref{Lagrangian contact structures.},
the fundamental GLA of contact type
$\mathfrak{m}= \mathfrak{g}_{-2}\oplus \mathfrak{g}_{-1}$ and
the decomposition
$\mathfrak{g}_{-1}=\mathfrak{e}^1\oplus \mathfrak{e}^2$ into the Lagrangian pair.
Moreover we f\/ix a volume form $\Omega_1 \in \wedge^n(\mathfrak{e}^{1*})$ on $\mathfrak{e}^1$ and
a volume form $\Omega_2 \in \wedge^n(\mathfrak{e}^{2*})$ on $\mathfrak{e}^2$.
Consider the group $C(\mathfrak{m}; \mathfrak{e}^1, \Omega_1; \mathfrak{e}^2, \Omega_2)$
consisting of all $a \in  {\rm GL}(\mathfrak{m})$ which satisf\/ies the following conditions:
$a\mathfrak{g}_{-1} = \mathfrak{g}_{-1}$,
the graded linear automorphism $\overline{a}$ of $\mathfrak{m}$ induced by $a$ is a GLA-automorphism,
$a\mathfrak{e}^1 = \mathfrak{e}^1$, $a\mathfrak{e}^2 = \mathfrak{e}^2$, and
$a^*\Omega_1 = \lambda\Omega_1$, $a^*\Omega_2 = \lambda\Omega_2$ for some $\lambda \in {\mathbf{R}}^{\times}$.
Then we have that $C(\mathfrak{m}; \mathfrak{e}^1, \Omega_1; \mathfrak{e}^2, \Omega_2)$ is identical with
$\overline{G}$.
\end{proof}

Thus the equivalence problem of bi-decomposable Monge--Amp{\`e}re systems with Lagrangian pairs
is studied as an adapted $G$-structure on a contact manifold of dimension~$2n+1$.

Let $G^0$ be a subgroup of ${\rm GL}(n+2, {\mathbf{R}})$ def\/ined by
\begin{gather*}
G^0 = \left\{ \left.
\begin{pmatrix}
                      \ k^{-1} & 0 & \ 0 \\
                      \ b_1 & A & \ 0 \\
                      \ a & {}^tb_2 &  \ k
                    \end{pmatrix}
                   \,  \right\vert \,
                    k \in {\mathbf{R}}^{\times}, \, A \in  {\rm SL}(n, {\mathbf{R}}), \, b_1, b_2 \in {\mathbf{R}}^n, \, a \in {\mathbf{R}}
                    \right\}.
\end{gather*}
We set
\begin{gather*}
H^0 = \left\{ \left.
\begin{pmatrix}
                      \ k^{-1} & 0 & \ 0 \\
                      \ 0 & A & \ 0 \\
                      \ 0 & 0 &  \ k
                              \end{pmatrix}
         \,  \right\vert   \,
        k \in {\mathbf{R}}^{\times}, \, A \in  {\rm SL} (n, {\mathbf{R}})
 \right\}
\subset G^0.
\end{gather*}
Then we have the model space
$G^0/H^0 \cong {\mathbf{R}}^{2n+1}$ with coordinates
$(x, z, p)$.
The $G^0$-action on ${\mathbf{R}}^{2n+1}$ is described by
\begin{gather*}
(x, z, p) \mapsto (x', z', p') = \left(k(Ax + b_1), k\big(kz - {}^tb_2x - a\big), k\left(p - \frac{1}{k}{}^tb_2\right)A^{-1}\right).
\end{gather*}
Then $dz' - p'dx' = k^2(dz - pdx)$ and
the form $\omega = c dx_1\wedge \cdots\wedge dx_n - dp_1\wedge\cdots\wedge dp_n$
is transformed to $\omega' = k^n\omega$.

Let $E_{i,j}$, $0 \leq i, j \leq n+1$, denotes the elementary $(n+2)\times (n+2)$-matrix such that
the $(i, j)$ component is $1$ and other components are all zero.
We set
\begin{gather*}
\varepsilon = - E_{0,0} + E_{n+1,n+1}, \qquad e_i = E_{i,0} \quad (1 \leq i \leq n), \\
f_j = E_{n+1,j} \quad (1 \leq j \leq n), \qquad \gamma = E_{n+1,0},
\end{gather*}
and identify the set of traceless matrices ${\mathfrak{sl}}(n, {\mathbf{R}})$ with
\begin{gather*}
\left\{ \left.
\begin{pmatrix}
                      0 & 0 &  0 \\
                       0 & A &  0 \\
                       0 & 0 &  0
                              \end{pmatrix}
         \,  \right\vert   \,
       A \in {\mathfrak{sl}}(n, {\mathbf{R}})
 \right\}
\subset {\mathfrak g}^0.
\end{gather*}
Then we set
${\mathfrak g}_0 = \langle\varepsilon\rangle_{{\mathbf{R}}} \oplus {\mathfrak{sl}}(n, {\mathbf{R}})$,
${\mathfrak g}_{-1}^1 = \langle e_1, \dots, e_n\rangle_{{\mathbf{R}}}$,
${\mathfrak g}_{-1}^2 = \langle f_1, \dots, f_n\rangle_{{\mathbf{R}}}$,
${\mathfrak g}_{-1} = {\mathfrak g}_{-1}^1 \oplus {\mathfrak g}_{-1}^2$,
and ${\mathfrak g}_{-2} = \langle\gamma\rangle_{{\mathbf{R}}}$.
Then
\begin{gather*}
{\mathfrak g}^0 = {\mathfrak g}_{-2} \oplus {\mathfrak g}_{-1} \oplus {\mathfrak g}_{0}
\end{gather*}
is the Lie algebra of $G^0$. We write
${\mathfrak m} = {\mathfrak g}_{-2}\oplus{\mathfrak g}_{-1}$.
Then the Lie subalgebra~${\mathfrak m}$ becomes the split Heisenberg algebra.

The matrix representation of $A_0 \in \mathfrak{g}_0=\langle\varepsilon\rangle_{{\mathbf{R}}} \oplus {\mathfrak{sl}}(n, {\mathbf{R}})$
with respect to the basis $\{ \gamma, \, e_i $ $(1 \leq i \leq n),\,   f_j\,  (1 \leq i \leq n) \}$ in
the Heisenberg algebra
$V= \mathfrak{m} = \mathfrak{g}_{-2}\oplus \mathfrak{g}_{-1}$
has the following form:
\begin{gather*}
A_0=C+A_0'=
 c
\begin{pmatrix}
2 & 0 & 0 \\
0 & I & \mathrm{O} \\
0 & \mathrm{O} & I
\end{pmatrix}
+
\begin{pmatrix}
0 & 0 & 0 \\
0 & A & \mathrm{O} \\
0 & \mathrm{O} & -{}^t\!A
\end{pmatrix}.
\end{gather*}
In fact, for the commutators, we have by the direct calculations:
\begin{gather*}
[\varepsilon, \gamma] = 2\gamma,
\qquad
[\varepsilon, e_i] = e_i,
\qquad
[\varepsilon, f_j] = f_j,
\\
[A, \gamma] = O,
\qquad
[A, e_i] = a_i,
\qquad
[A, f_j] = - a^j,
\end{gather*}
where $a_i$ is the $i$-th column of $A$ and $a^j$ is the $j$-th row of $A$,
$1 \leq i, j \leq n$, and $A \in {\mathfrak{sl}}(n, {\mathbf{R}})$.
Moreover we have
\begin{gather*}
[e_i, f_j] = - \delta_{ij}\gamma, \qquad [e_i, e_j] = 0, \qquad [f_i, f_j] = 0 \quad (1 \leq i, j \leq n).
\end{gather*}

We will study the prolongation of
$({\mathfrak m}, {\mathfrak g}_0)$.
We def\/ine the prolongation inductively
${\mathfrak g}_k = {\mathfrak g}({\mathfrak m}, {\mathfrak g}_0)_k$ $(k \geq 1)$ by
the set of elements
$\{
(\alpha, \beta) \in  {\rm Hom}({\mathfrak g}_{-1}, {\mathfrak g}_{k-1})\oplus
{\rm Hom}({\mathfrak g}_{-2}, {\mathfrak g}_{k-2})\}$ satisfying
\begin{alignat*}{3}
& {\rm (i)} \ \ && \beta([x, y]) = [\alpha(x), y] - [\alpha(y), x]   \qquad (x, y \in {\mathfrak g}_{-1}), &
\\
& {\rm (ii)} \ \ && [\alpha(y), z] = [\beta(z), y]  \qquad (y \in {\mathfrak g}_{-1},\, z \in {\mathfrak g}_{-2}).&
\end{alignat*}
See \cite[p.~429]{Y}. Then we have

\begin{Lem}
\label{prolongation-calculation}
The prolongation ${\mathfrak g}_i$ vanishes for any $i \geq 1$.
\end{Lem}

\begin{proof}
First we calculate ${\mathfrak g}_1$. For any $(\alpha, \beta) \in
{\rm Hom}(\mathfrak{g}_{-1}, \mathfrak{g}_0) \oplus {\rm Hom}(\mathfrak{g}_{-2}, \mathfrak{g}_{-1})$,
the condi\-tions~(i) and (ii) imply that
\begin{alignat*}{3}
& \text{(i)}   \ \ && \beta([e_i, f_j]) = [\alpha(e_i), f_j] - [\alpha(f_j), e_i]  \qquad (1 \leq i, j \leq n), &
\\
& \text{(ii-1)} \ \ & & [\alpha(e_i), \gamma] = [\beta(\gamma), e_i] \qquad (1 \leq i \leq n),&
\\
& \text{(ii-2)} \ \ && [\alpha(f_j), \gamma] = [\beta(\gamma), f_j] \qquad (1 \leq j \leq n).&
\end{alignat*}
Set
\begin{gather*}
\beta(\gamma) = \sum_{\ell=1}^n  b_\ell e_\ell + \sum_{m=1}^n  c_m  f_m,
\qquad
\alpha(e_i) = h_i\varepsilon + A^{(i)}, \qquad
\alpha(f_j) = k_j\varepsilon + B^{(j)},
\end{gather*}
where $b_\ell, c_m, h_i, k_j \in {\mathbf{R}}$, $A^{(i)}, B^{(j)} \in {\mathfrak{sl}}(n, {\mathbf{R}})$.
Then
\begin{gather*}
\beta([e_i, f_j]) = \sum_{\ell=1}^n (- \delta_{ij}b_\ell)e_\ell + \sum_{m=1}^n (- \delta_{ij}c_m)f_m,
\\
{[\alpha(e_i), f_j]} = h_i f_j - \big(j{\mbox{\rm -th row of }} A^{(i)}\big), \qquad
{[\alpha(f_j), e_i]} = k_j e_i + \big(i{\mbox{\rm -th column of }} B^{(j)}\big),
\\
{[\alpha(e_i), \gamma]} = 2h_i\gamma, \qquad {[\beta(\gamma), e_i]} = c_i\gamma, \qquad
{[\alpha(f_j), \gamma]} = 2k_j\gamma, \qquad {[\beta(\gamma), f_j]} = - b_j\gamma.
\end{gather*}
By the condition (i), we have
\begin{gather*}
B^{(j)}_{\ell i} = - \delta_{\ell i}k_j + \delta_{ij}b_\ell, \qquad
A^{(i)}_{j m} = \delta_{j m}h_i + \delta_{ij}c_m \qquad (1 \leq i, j, \ell, m \leq n).
\end{gather*}
By the condition (ii-1), we have $c_i = 2h_i$, $1 \leq i \leq n$.
By the condition (ii-2), we have $b_j = - 2k_j$, $1 \leq j \leq n$.
Therefore we have
\begin{gather*}
B^{(j)}_{\ell i} = - \delta_{\ell i}k_j - 2\delta_{ij}k_\ell, \qquad
A^{(i)}_{j m} = \delta_{j m}h_i + 2\delta_{ij}h_m \qquad (1 \leq i, j, \ell, m \leq n).
\end{gather*}
In particular, for any $j$, $\ell$ with $1 \leq j, \ell \leq n$,
we have $B^{(j)}_{\ell\ell} = - k_j - 2\delta_{\ell j}k_\ell$, and therefore
\begin{gather*}
0 = \operatorname{tr} B^{(j)} = - (n+2)k_j,
\end{gather*}
hence $k_j = 0$, $1 \leq j \leq n$.
For any $i$, $m$ with $1 \leq i, m \leq n$, we have
$A^{(j)}_{mm} = h_i + 2\delta_{i m}h_m$, and therefore
\begin{gather*}
0 = \operatorname{tr}  A^{(i)} = (n+2)h_i,
\end{gather*}
hence $h_i = 0$, $1 \leq i \leq n$.
Thus we have $\beta = 0$ and $\alpha = 0$. Therefore we have ${\mathfrak g}_1 = 0$.

Second we calculate ${\mathfrak g}_2$. Take any $(\alpha, \beta) \in {\mathfrak g}_2 \subset
{\rm Hom}({\mathfrak g}_{-1}, {\mathfrak g}_{1})\oplus
{\rm Hom}({\mathfrak g}_{-2}, {\mathfrak g}_{0})$.
Since ${\mathfrak g}_{1} = 0$, $\alpha = 0$. Then we see $\beta(\gamma) = 0$. Therefore $\beta = 0$.
Thus we obtain that ${\mathfrak g}_{2} = 0$.

From ${\mathfrak g}_{1} = 0$, ${\mathfrak g}_{2} = 0$, we have ${\mathfrak g}_{i} = 0$, $i \geq 3$, automatically.
\end{proof}

Then we have

\begin{Th}
\label{maximal symmetry}
Let $(M, D)$ be a contact manifold of dimension $2n+1$
and~$\mathcal{M}$
a bi-decomposable Monge--Amp{\`e}re system
with a Lagrangian pair on~$(M,D)$.
Assume that $n \geq 3$.
Then the automorphism pseudo-group $\operatorname{Aut}(\mathcal{M})$
of~$\mathcal{M}$ has
at most the dimension of
$(n+1)^2$. The estimate is best possible.
\end{Th}

\begin{proof}
By Lemma \ref{prolongation-calculation},
we have known that the prolongations ${\mathfrak g}_i$, $1 \leq i$ in the sense of Tanaka vanish.
On the other hand $\overline{G}$ is equal to $(H^0)^{\#}$ in the notation of~\cite{Tan}.
Then, by the f\/initeness theorem of Tanaka~\cite[Corollary~2]{Tan}, we see that
the dimension of the pseudo-group of automorphisms on ${\mathcal M}$ is estimated by
$\dim({\mathfrak g^0}) = \sum\limits_{i \leq 0} \dim({\mathfrak g}_i) = (n+1)^2$.
Moreover there exists an Monge--Amp{\`e}re system $\mathcal{M}$ with a Lagrangian pair on
$M = {\mathbf{R}}^{2n+1}$, which arises in equi-af\/f\/ine geometry, such that
the automorphism group {\rm Aut}($\mathcal{M}$) attains the maximal dimension $(n+1)^2$.
See Section~\ref{Homogeneous M-A systems with Lagrangian pairs.}.1.
\end{proof}
